\section{Smoothness and a proper intrinsic distance}\label{sec:geometry}

We first establish two consequences of boundedness and cocompactness that do
not use semialgebraicity: the open set is smooth, and analytic discs cannot
escape it during a deformation whose boundaries remain inside.  These facts
will allow us to work at the boundary without assuming its regularity.
The peak-function localization has its classical predecessor in Rosay's
argument at a strictly pseudoconvex boundary point
\cite[Section~3, Lemma~2]{Rosay1979}. Here a finite-projection argument
allows the localization to begin before the open set is known to be smooth.

Embed the initial reduced affine variety in $\C^N$ in coordinates in which
$U$ is bounded.  Its closure $\overline U$ is compact.  There is a compact set
$K\subset U$ such that $\Gamma K=U$.  Indeed, the quotient map is open, and
the images of relatively compact open neighborhoods in $U$ cover $U/\Gamma$;
finitely many suffice.  Their closures give $K$.  This argument does not
require the action to be free.

\begin{lemma}\label{lem:singular-montel}
Let $W$ be an open subset of a reduced complex affine algebraic set.  The
family of holomorphic functions $h$ on $W$ with $|h|\leq 1$ is locally
equicontinuous.  Every sequence in this family has a subsequence converging
uniformly on compact subsets to a continuous function $h_\infty$.
If $h_\infty(q)=1$ at a point $q\in W$, then $h_\infty=1$ in a neighborhood
of $q$.  In particular, if $W$ is connected, then $h_\infty\equiv 1$.
\end{lemma}

\begin{proof}
We give the argument at a possibly singular and nonnormal point $q$.  Work
first on one irreducible algebraic component $V_0$ through $q$, of positive
dimension $d$.  Noether normalization
\cite[Lemma~10.115.4, Tag~00OY]{StacksNoether} gives a finite algebraic map
$\pi:V_0\to\C^d$.  Translate so that $\pi(q)=0$.  By properness and
finiteness, we may restrict over a sufficiently small polydisc $B$ about $0$
and retain just the part $\Omega$ near $q$, with
\[
  \pi:\Omega\longrightarrow B\quad\hbox{proper and finite},
  \qquad \pi^{-1}(0)\cap\Omega=\{q\},\qquad \Omega\subset W.
\]
To make this restriction, take disjoint neighborhoods of the finitely many
points over $0$ and shrink $B$ until its entire inverse image lies in their
union.  The part near $q$ is then both open and closed in that inverse image.
Outside a proper algebraic hypersurface $E$ in $B$, the restricted map is a
covering with a fixed number $r\geq1$ of holomorphic sheets.  This follows
also directly from separability of the finite function-field extension:
discard the denominators and discriminant of a primitive element.

For $|h|\leq1$, form the monic polynomial of its sheet values,
\[
  P_h(z,T)=\prod_{x\in\pi^{-1}(z)\cap\Omega}(T-h(x))
           =T^r+\sum_{\nu=0}^{r-1}a_{\nu,h}(z)T^\nu,
  \qquad z\in B\setminus E.
\]
The coefficients are single-valued holomorphic functions bounded by
constants depending only on $r$.  They extend holomorphically across $E$
by the removable-singularity theorem for bounded holomorphic functions.
One elementary proof uses a complex line through the point that is not
contained in $E$, a small circle on that line missing $E$, and the Cauchy
integral on nearby parallel circles.  Properness and continuity of $h$ give
\[
                         P_h(0,T)=(T-h(q))^r.
\]
The polynomial identity $P_h(\pi(x),h(x))=0$ holds on all of $\Omega$ by
continuity, since the sheet locus is dense.  Cauchy estimates for the
uniformly bounded coefficients, on a smaller polydisc, therefore yield
\[
  |h(x)-h(q)|^r
   =|P_h(0,h(x))-P_h(\pi(x),h(x))|
   \leq C\|\pi(x)\|.
\]
Here $C$ is independent of $h$.  This proves equicontinuity at $q$ on $V_0$.
There are only finitely many components through $q$; applying the argument
to each proves the assertion on their union.  A zero-dimensional component
requires no argument.  Local compactness, a countable compact exhaustion,
and the Arzel\`a--Ascoli theorem now give compact convergence of a subsequence.

For the maximum assertion, use the same chart and consider the sheet
average
\[
  b_h(z)=\frac1r\sum_{x\in\pi^{-1}(z)\cap\Omega}h(x).
\]
It extends holomorphically to $B$, satisfies $|b_h|\leq1$, and has
$b_h(0)=h(q)$.  If $h_j\to h_\infty$ locally uniformly and
$h_\infty(q)=1$, ordinary normal-family compactness on $B$ gives, after
another extraction, a holomorphic limit of the sheet averages with value
$1$ at $0$.  The usual
maximum principle makes that limit identically $1$.  On every regular fiber,
an average of numbers in the closed unit disc can equal $1$ only if every
number equals $1$.  Thus $h_\infty=1$ on the sheet locus near $q$, hence
everywhere near $q$ by continuity.  Apply this on every component through
$q$.  Finally, $\{h_\infty=1\}$ is both closed and open in $W$, which proves
the connected case.  No identification of weakly holomorphic functions
with holomorphic functions on a nonnormal space was used.
\end{proof}

\begin{proposition}\label{prop:smoothness}
If $U$ is not a point, then it is a complex manifold of positive dimension
and lies in the smooth locus of a single irreducible component $V$ of the
ambient variety.
\end{proposition}

\begin{proof}
Choose a positive-dimensional irreducible component $V_0$ whose smooth
points not shared with another component meet $U$.  Such a component exists:
these points are dense in each positive-dimensional component, whereas a
zero-dimensional irreducible component is isolated and cannot meet the
connected nonsingleton $U$.  There is a polynomial $P$ vanishing on the
other components and on the singular locus of $V_0$, but not identically on
$V_0$.  In particular, $P$ is nonzero somewhere in $U$.

Maximize $\log|P(q)|+\|q\|^2$ on $\overline U$, with value $-\infty$ at
zeros of $P$.  A maximizer $p$ has $P(p)\ne0$, so it is a smooth, unshared
point of $V_0$.  It belongs to $\partial U$: near such a point the function
is strictly plurisubharmonic, and hence cannot have an interior maximum.
The entire function
\[
  h(q)=\frac{P(q)}{P(p)}
          \exp\bigl(2\overline p\cdot(q-p)\bigr)
\]
satisfies
\begin{equation}\label{eq:initial-peak}
  |h(q)|\leq \exp\bigl(-\|q-p\|^2\bigr)
  \quad(q\in\overline U),\qquad h(p)=1.
\end{equation}
Indeed, the logarithm of the left-hand side is at most
$\|p\|^2-\|q\|^2+2\operatorname{Re}(\overline p\cdot(q-p))$.

Choose $q_j\in U$ tending to $p$, and write $q_j=\gamma_j k_j$ with
$k_j\in K$.  After extraction, $k_j\to k\in K$, and
Lemma~\ref{lem:singular-montel} gives locally uniform convergence of
$h\circ\gamma_j$.  Equicontinuity at the moving points $k_j$ shows that
the limit has value $1$ at $k$.  It is therefore identically $1$ on $U$.
Equation~\eqref{eq:initial-peak} implies that $\gamma_j$ converges to $p$
uniformly on every compact subset of $U$.

For each $q\in U$, the point $\gamma_j(q)$ eventually lies in a smooth
ambient neighborhood of $p$.  Since a biholomorphism preserves smoothness
of the complex space, $q$ is smooth.  A smooth point belongs to exactly one
irreducible component of the ambient variety; the component is locally
constant on $U$, and hence constant by connectedness.  This gives $V$.
\end{proof}

The singleton is the zero-dimensional bounded symmetric domain, so henceforth
we work on this irreducible $V$ and put $n=\dim_\C U\geq1$.

We use Kobayashi's intrinsic distance, defined through holomorphic disc
chains~\cite[Section~2]{Kobayashi1967}. Its contraction property is built
into this definition; properness here will follow from cocompactness.
Let $\rho$ denote the Poincar\'e distance on the unit disc $\D$.  For $x,y\in U$,
define $d_U(x,y)$ as the infimum of
\[
                    \sum_{j=1}^s\rho(\alpha_j,\beta_j)
\]
over finite chains of holomorphic maps $a_j:\D\to U$ with
$a_1(\alpha_1)=x$, $a_s(\beta_s)=y$, and
$a_j(\beta_j)=a_{j+1}(\alpha_{j+1})$.  By its definition, holomorphic maps
decrease this chain-of-discs distance, and biholomorphisms preserve it.

\begin{proposition}\label{prop:proper-distance}
The function $d_U$ is a finite distance inducing the usual topology of $U$.
Every closed $d_U$-ball is compact.
\end{proposition}

\begin{proof}
Coordinate balls give finite disc chains locally, hence globally by
connectedness.  Bounded ambient coordinate functions and the Schwarz lemma
give a constant $c>0$ such that
\begin{equation}\label{eq:distance-lower-bound}
                         d_U(x,y)\geq c\|x-y\|.
\end{equation}
For example, send each bounded coordinate to $\D$; the Poincar\'e distance
there dominates a constant times Euclidean distance, and sum along a chain.
Conversely, small affine complex discs in a manifold chart give a local
upper bound tending to zero with the coordinate separation.  Thus $d_U$ is
a distance and induces the original topology.

Choose a relatively compact neighborhood of $K$ in $U$.
Equation~\eqref{eq:distance-lower-bound} and compactness of $K$ give an
$r>0$ for which every closed $r$-ball centered in $K$ lies in that
neighborhood's compact closure.  Such a ball is closed in the original
topology, hence compact.  Translating by $\Gamma$ shows that every closed
$r$-ball in $U$ is compact, with the same $r$.

We spell out why this uniform local statement implies properness.  Fix $x$,
and suppose $\overline B_{d_U}(x,R)$ is compact.  Cover it by finitely many
open $r/4$-balls centered at $x_1,\ldots,x_l$.  For
$d_U(x,y)\leq R+r/4$, choose a disc chain of total Poincar\'e length less than
$d_U(x,y)+r/4$.  Replace each disc link by a Poincar\'e geodesic segment.
If the total length exceeds $R$, the point at accumulated length $R$
lies in $\overline B_{d_U}(x,R)$, and the remaining length is less than
$r/2$.  Hence $y$ belongs to an $r$-ball centered at some $x_i$.
If the total length is at most $R$, the same conclusion is immediate.
Consequently
\[
  \overline B_{d_U}(x,R+r/4)
       \subset\bigcup_{i=1}^l\overline B_{d_U}(x_i,r).
\]
The set on the left is closed and the union on the right is compact.
Starting at $R=r$ and iterating proves compactness for every radius.
\end{proof}

\begin{lemma}[Continuity principle]\label{lem:continuity}
Let $\Omega$ be a smooth ambient coordinate neighborhood, and let
$A:[0,1]\times\overline\D\to\Omega$ be continuous, with each
$A_t=A(t,\cdot)$ holomorphic on $\D$.  Suppose
\[
  A_t(\partial\D)\subset U\quad(0\leq t\leq1),
  \qquad A_0(\overline\D)\subset U.
\]
Then $A_t(\overline\D)\subset U$ for every $t$.  In particular, this applies
to continuous deformations of closed $C^1$ analytic discs; holomorphic
extension through the boundary circle is not required.
\end{lemma}

\begin{proof}
The set $T=\{t:A_t(\overline\D)\subset U\}$ is relatively open in $[0,1]$
by compactness of the closed disc.  To see that it is closed, let
$t_j\in T$ tend to $t_*$.  The boundary condition and continuity give a
fixed collar $1-\eps\leq|\zeta|\leq1$ whose images, for $t$ near $t_*$,
lie in a compact set $L\subset U$.  Fix $\zeta_0$ in the interior of this
collar and $x_0\in U$.  By Proposition~\ref{prop:proper-distance},
$M=\max_{x\in L}d_U(x_0,x)$ is finite.  For each fixed $\zeta\in\D$,
the earlier discs take values in $U$, so Schwarz contraction gives
\[
  d_U(x_0,A_{t_j}(\zeta))
            \leq M+\rho(\zeta_0,\zeta).
\]
Properness places these values in a fixed compact subset of $U$.
Their ambient limit $A_{t_*}(\zeta)$ therefore belongs to $U$.
The boundary already lies in $U$, so $t_*\in T$.  Since $T$ is nonempty,
open, and closed in $[0,1]$, it is all of $[0,1]$.
\end{proof}
